\documentclass[12pt,a4paper]{report}

% ==========================================
% PACKAGES & MARGINS
% ==========================================
\usepackage[utf8]{inputenc}
\usepackage[T1]{fontenc}
\usepackage{times} % Standard Times font
\usepackage[top=1in,bottom=1in,left=1.25in,right=1in]{geometry}
\usepackage{setspace}
\onehalfspacing

\usepackage{graphicx}
\usepackage{amsmath,amssymb,amsfonts}
\usepackage{enumitem}
\usepackage{booktabs}
\usepackage{multirow}
\usepackage{array}
\usepackage{tabularx}
\usepackage{caption}
\captionsetup[table]{font=small,labelfont=bf,skip=8pt}
\captionsetup[figure]{font=small,labelfont=bf,skip=8pt}
\usepackage{hyperref}
\hypersetup{
    colorlinks=true,
    linkcolor=black,
    citecolor=black,
    urlcolor=blue
}
\usepackage{cite}
\usepackage{titlesec}
\usepackage{fancyhdr}

% Chapter heading formatting
\titleformat{\chapter}[display]
  {\normalfont\Large\bfseries\centering}{\chaptertitlename\ \thechapter}{10pt}{\LARGE}
\titlespacing*{\chapter}{0pt}{-20pt}{25pt}

% Section formatting
\titleformat{\section}
  {\normalfont\large\bfseries}{\thesection}{1em}{}
\titleformat{\subsection}
  {\normalfont\normalsize\bfseries}{\thesubsection}{1em}{}

% Header and Footer
\pagestyle{fancy}
\fancyhf{}
\cfoot{\thepage}
\renewcommand{\headrulewidth}{0pt}

% Typography & Table Spacing
\raggedbottom
\renewcommand{\arraystretch}{1.25}
\newcolumntype{C}{>{\centering\arraybackslash}X}
\newcolumntype{L}{>{\raggedright\arraybackslash}X}

\begin{document}

% ==========================================
% 1. FRONT MATTER
% ==========================================
\begin{titlepage}
    \centering
    \vspace*{0.5cm}
    {\large \textbf{A PROJECT REPORT ON}}\\[0.8cm]
    {\LARGE \textbf{MEASURING AND REMEDIATING SEMANTIC CONSOLIDATION DRIFT IN LONG-HORIZON LLM AGENTS}}\\[1.5cm]
    
    {\large \textit{Submitted in partial fulfillment of the requirements for the degree of}}\\[0.3cm]
    {\large \textbf{BACHELOR OF TECHNOLOGY}}\\[0.2cm]
    {\large \textit{in}}\\[0.2cm]
    {\large \textbf{COMPUTER SCIENCE AND ENGINEERING}}\\[1.2cm]
    
    {\large \textit{by}}\\[0.4cm]
    \textbf{AARAV RAJ SHRESTHA} \quad (Reg. No: \textbf{AP23110011302})\\[1.2cm]
    
    {\large \textit{Under the guidance of}}\\[0.3cm]
    \textbf{[Guide Name / Faculty Mentor]}\\[0.2cm]
    {Assistant Professor, Department of Computer Science and Engineering}\\[1.5cm]
    
    \vspace*{0.5cm}
    {\large \textbf{DEPARTMENT OF COMPUTER SCIENCE AND ENGINEERING}}\\[0.2cm]
    {\large \textbf{SRM UNIVERSITY-AP, ANDHRA PRADESH}}\\[0.2cm]
    {\large \textbf{DECEMBER 2026}}
\end{titlepage}

\pagenumbering{roman}
\setcounter{page}{2}

% ------------------------------------------
% DECLARATION
% ------------------------------------------
\chapter*{DECLARATION}
\addcontentsline{toc}{chapter}{DECLARATION}

I hereby declare that this project report entitled \textbf{``Measuring and Remediating Semantic Consolidation Drift in Long-Horizon LLM Agents''} submitted in partial fulfillment of the requirements for the degree of \textbf{Bachelor of Technology in Computer Science and Engineering} at \textbf{SRM University-AP}, is an authentic record of my own work carried out under the supervision of my project guide.

To the best of my knowledge, the matter embodied in this report has not been submitted to any other University or Institute for the award of any degree or diploma. I further declare that the intellectual content of this report is the product of my own research and that all assistance received in preparing the report and sources utilized have been duly acknowledged in IEEE format. The similarity index and AI-generated text thresholds comply with university norms ($<20\%$).

\vspace{2cm}
\noindent
\textbf{Place:} Amaravati, Andhra Pradesh\\
\textbf{Date:} \today

\vspace{1.5cm}
\begin{flushright}
\textbf{Aarav Raj Shrestha}\\
(Reg. No: AP23110011302)\\
Department of Computer Science and Engineering\\
SRM University-AP
\end{flushright}

\clearpage

% ------------------------------------------
% CERTIFICATE
% ------------------------------------------
\chapter*{CERTIFICATE}
\addcontentsline{toc}{chapter}{CERTIFICATE}

This is to certify that the project report entitled \textbf{``Measuring and Remediating Semantic Consolidation Drift in Long-Horizon LLM Agents''} is a bonafide work carried out by \textbf{Aarav Raj Shrestha (Reg. No: AP23110011302)} in partial fulfillment of the requirements for the award of the degree of \textbf{Bachelor of Technology in Computer Science and Engineering} from \textbf{SRM University-AP, Andhra Pradesh}, during the academic period 2026 under my supervision and guidance.

\vspace{3cm}
\noindent
\begin{tabular}{p{7.5cm}p{7.5cm}}
\textbf{[Guide Name]} & \textbf{Dr. V M Manikandan} \\
Project Guide & Head of the Department \\
Assistant Professor & Associate Professor \& Head \\
Department of CSE & Department of CSE \\
SRM University-AP & SRM University-AP
\end{tabular}

\clearpage

% ------------------------------------------
% ACKNOWLEDGMENT
% ------------------------------------------
\chapter*{ACKNOWLEDGMENT}
\addcontentsline{toc}{chapter}{ACKNOWLEDGMENT}

First and foremost, I would like to express my profound gratitude and deep reverence to my project supervisor, \textbf{[Guide Name]}, for invaluable guidance, patient mentoring, rigorous critique, and constant encouragement throughout the course of this Undergraduate Research Opportunity Project (UROP). 

I convey my heartfelt gratitude to \textbf{Dr. V M Manikandan}, Head of the Department of Computer Science and Engineering, for providing excellent academic infrastructure, computational facilities, and continuous departmental encouragement to conduct quality research.

I extend my thanks to the entire faculty and technical staff of the Department of Computer Science and Engineering at SRM University-AP for their direct and indirect support. Lastly, I express my sincere appreciation to my peers, group members, and family, whose unfailing moral support, insightful discussions, and patience sustained my efforts during experimental executions and report writing.

\vspace{1.5cm}
\begin{flushright}
\textbf{Aarav Raj Shrestha}\\
(AP23110011302)
\end{flushright}

\clearpage

% ------------------------------------------
% ABSTRACT
% ------------------------------------------
\chapter*{ABSTRACT}
\addcontentsline{toc}{chapter}{ABSTRACT}

Autonomous Large Language Model (LLM) agents operating over extended conversational horizons must periodically compress, consolidate, and summarize historical interactions into persistent working memory due to context window constraints. Existing memory benchmarks evaluate retention almost exclusively through binary fact recall and retrieval fidelity. This investigation examines whether iterative memory consolidation acts merely as a stochastic information-loss channel ($\mathcal{H}_0$) or induces systematic, directional semantic drift ($\mathcal{H}_1$). We formulate the DriftBench evaluation suite, establishing fine-grained 5-point semantic scales across deontic obligation, preference intensity, epistemic certainty, numerical magnitude, and polarity negation.

Through rigorous methodological remediation—resolving token exhaustion bugs, implementing line-level chunking, and establishing an item-clustered block bootstrap to correct for repeated-measures autocorrelation—we benchmark sliding window, hierarchical, and Anchor-Verify-Restore (AVR) memory architectures across 5,760 evaluated states. Our pre-registered Stage A experiment ($N=48$ items, 20 rounds, multi-turn contexts) demonstrates that while attractor collapse toward neutral defaults ($\beta_{\text{ARC}} \to -1.0$) is definitively rejected as an experimental artifact, hierarchical consolidation induces statistically significant directional strengthening on obligation and preference attributes (+0.33 scale points, $p=0.0028$). Control attributes remain invariant ($-0.03$ points). To mitigate memory decay, we deploy AVR, which slashes contextual attrition from 21.1\% to 1.1\% and elevates exact ground-truth retention from 78.0\% to 87.4\% ($p=0.047$). These findings confirm that agent memory summarization is fundamentally biased toward categorical escalation and establish algorithmic foundations for persistent, verifiably anchored agent state architectures.

\clearpage

% ------------------------------------------
% CONTENTS, LIST OF TABLES, LIST OF FIGURES
% ------------------------------------------
\tableofcontents
\clearpage

\listoftables
\addcontentsline{toc}{chapter}{LIST OF TABLES}
\clearpage

\listoffigures
\addcontentsline{toc}{chapter}{LIST OF FIGURES}
\clearpage

% ==========================================
% 2. MAIN BODY CHAPTERS
% ==========================================
\pagenumbering{arabic}
\setcounter{page}{1}

% ------------------------------------------
% CHAPTER 1: INTRODUCTION
% ------------------------------------------
\chapter{Introduction to the Project}

\section{Background and Domain Overview}
Recent advances in generative artificial intelligence have enabled the deployment of autonomous Large Language Model (LLM) agents capable of multi-step planning, long-horizon tool execution, and continuous conversational interaction. Real-world agent deployments—spanning personal productivity co-pilots, enterprise customer support systems, and autonomous research assistants—interact with human users over weeks or months. However, Transformer-based language models are fundamentally constrained by finite context windows and high computational attention complexity ($O(N^2)$). To sustain coherent operation beyond token boundaries, agents must compress historical dialogue turns into condensed, persistent working memory.

\section{Problem Statement}
Existing literature on conversational memory (e.g., MemGPT, LongMemEval) evaluates memory systems almost exclusively through binary information recall: whether a key fact can be retrieved or whether it has been lost. This paradigm overlooks a subtle, potentially hazardous phenomenon: \textbf{semantic consolidation drift}. When an LLM iteratively summarizes previous summaries, does memory degradation manifest purely as omission and information loss, or does the summarization process introduce directional distortion in meaning? Specifically, do soft suggestions (e.g., ``I should water the plants'') escalate into strict mandates (``I must water the plants''), or do nuanced preferences collapse toward generic defaults?

\section{Project Objectives}
The core objectives of this Undergraduate Research Opportunity Project (UROP) are:
\begin{itemize}
    \item \textbf{Develop DriftBench v1:} Construct a standardized, fine-grained benchmark dataset containing 200 calibrated statements mapped to 5-point semantic intensity scales across deontic obligation, preference strength, epistemic certainty, numerical magnitude, and polarity negation.
    \item \textbf{Formulate Robust Quantitative Metrics:} Define the Attractor Regression Coefficient (ARC), mean signed semantic drift, attrition rate, and item-clustered block bootstrap inference methods to account for repeated-measures autocorrelation across Markovian memory rounds.
    \item \textbf{Analyze Memory Architectures:} Empirically evaluate multi-round consolidation under Sliding Window and Hierarchical Summarization pipelines across varied contextual noise levels ($c0, c8, c20$).
    \item \textbf{Design and Validate Remediation (AVR):} Engineer an autonomous Anchor-Verify-Restore (AVR) framework that identifies critical semantic propositions, verifies them across consolidation rounds, and verbatim-restores corrupted or omitted anchors.
\end{itemize}

\section{Report Organization}
The remainder of this report is structured as follows: Chapter 2 details the motivation and practical urgency of studying memory drift. Chapter 3 provides a comprehensive literature survey contrasting existing memory architectures with our work. Chapter 4 elaborates the design, formal mathematical metrics, and benchmark architecture. Chapter 5 covers system implementation, experimental protocols, and bug remediations. Chapter 6 outlines the software and computational tools employed. Chapter 7 presents empirical results, statistical significance tests, and discussions. Finally, Chapter 8 concludes the project and charts concrete avenues for future research.

% ------------------------------------------
% CHAPTER 2: MOTIVATION
% ------------------------------------------
\chapter{Motivation}

\section{The Fragility of Agentic State}
Autonomous software agents increasingly handle high-stakes human responsibilities, including managing personal schedules, enforcing compliance protocols, and configuring automated enterprise workflows. In these domains, the precise modal nuance of user instructions is paramount. For example:
\begin{itemize}
    \item \textbf{Deontic Obligation:} A user stating, \textit{``I should submit the expense report this week, but it is not urgent,''} conveys a low-pressure priority. If an agent's internal memory progressively consolidates this into \textit{``The expense report is strictly due and must be submitted immediately,''} the agent may trigger premature reminders or initiate unwarranted transactions.
    \item \textbf{Preference Intensity:} A user stating, \textit{``I slightly prefer morning meetings, but afternoon is fine,''} should not have afternoon collaboration invites categorically rejected due to consolidated memory storing \textit{``Requires morning meetings.''}
    \item \textbf{Epistemic Certainty and Negation:} A tentative hypothesis (``The deployment might succeed by Friday'') or an explicit negative constraint (``Never sign without review'') can trigger catastrophic system failures if consolidated into unwarranted certainty or dropped negation.
\end{itemize}

\section{Limitations in Existing Memory Paradigms}
Existing agent memory frameworks operate under the assumption that LLM summarization is unbiased and information-preserving up to random loss. Consequently, memory benchmarks evaluate retrieval accuracy rather than semantic stability over successive compression passes. When summarizers shorten text, inductive biases favor concise, categorical assertions over qualified, hedged statements. Understanding whether these systematic distortions accumulate over successive consolidation cycles is critical for building safe, dependable, long-horizon AI assistants.

% ------------------------------------------
% CHAPTER 3: LITERATURE SURVEY
% ------------------------------------------
\chapter{Literature Survey}

\section{Evolution of LLM Memory Architectures}
Long-context reasoning has progressed along two primary axes: architectural context window expansion and externalized memory architectures. While models like Gemini 1.5 and Claude 3.5 process hundreds of thousands of tokens natively, infinite-context windows suffer from the ``lost-in-the-middle'' phenomenon and escalating inference costs. Consequently, external hierarchical memory systems remain indispensable for persistent agents.

Packer et al. \cite{packer2023memgpt} introduced \textbf{MemGPT}, an operating-system-inspired architecture that manages hierarchical context tiers (main context, working context, and archival storage) via explicit function calls. While MemGPT demonstrates impressive capabilities in multi-session chat, it treats memory writing as an unmonitored text-generation action, leaving the semantic fidelity of recurring edits unmeasured.

\section{Benchmark Datasets for Memory Evaluation}
Wu et al. \cite{wu2024longmemeval} established \textbf{LongMemEval}, a rigorous benchmark testing five core memory capabilities over long conversational histories: information extraction, multi-session reasoning, temporal reasoning, knowledge updating, and preference retention. Although LongMemEval evaluates preference recall, its evaluation framework is binary (fact present vs. absent) and does not track continuous shift along scalar dimensions over dozens of iterative summaries.

\section{Iterative Summarization and Information Decay}
Zhang et al. \cite{zhang2024cumulative} investigated cumulative error propagation in recursive summarization, showing that factual inaccuracies compound exponentially across generation chains. However, their work focused on document summarization rather than conversational agent states, and did not examine scalar modal shifts (obligation, certainty) or propose algorithmic restoration wrappers. Table \ref{tab:lit_survey} summarizes the research gap addressed by our work.

\begin{table}[htbp]
\centering
\small
\caption{Comparison of Agent Memory Paradigms and Benchmarks}
\label{tab:lit_survey}
\begin{tabularx}{\textwidth}{p{3.2cm}p{2.8cm}p{3.2cm}X}
\toprule
\textbf{Framework / Work} & \textbf{Memory Model} & \textbf{Evaluation Focus} & \textbf{Limitation Addressed in UROP} \\
\midrule
MemGPT \cite{packer2023memgpt} & OS-style tiered paging & Task completion, tool use & No continuous measurement of semantic drift across rewrites. \\
LongMemEval \cite{wu2024longmemeval} & Multi-session dialogues & Binary retrieval, fact recall & Lacks fine-grained 5-point semantic intensity scales and repeated consolidation tracking. \\
Cumulative Error \cite{zhang2024cumulative} & Recursive document summaries & Factual hallucination & Restricted to news documents; neglects agentic state and deontic modalities. \\
\textbf{DriftBench (This Work)} & \textbf{Sliding, Hierarchical, AVR} & \textbf{Directional drift, ARC, attrition, cluster-bootstrap} & \textbf{Provides continuous scalar calibration, pre-registered hypothesis testing, and active AVR repair.} \\
\bottomrule
\end{tabularx}
\end{table}

% ------------------------------------------
% CHAPTER 4: DESIGN AND METHODOLOGY
% ------------------------------------------
\chapter{Design and Methodology}

\section{System Architecture}
The experimental workflow consists of four modular stages: Benchmark Synthesis, Contextual Dialogue Embedding, Multi-Round Memory Consolidation, and Independent LLM-Based Evaluation. Figure \ref{fig:sys_arch} illustrates the complete experimental architecture.

\begin{figure}[htbp]
\centering
\fbox{\parbox{0.9\textwidth}{\centering
\vspace{0.5cm}
\textbf{[DriftBench v1: 200 Standardized Seed Items (Levels 1--5)]}\\[0.3cm]
$\downarrow$\\[0.3cm]
\textbf{[Contextual Dialogue Injection: $c0$ (0 turns), $c8$ (8 turns), $c20$ (20 turns)]}\\[0.3cm]
$\downarrow$\\[0.3cm]
\textbf{[Consolidation Engines: Sliding Window \quad vs. \quad Hierarchical \quad vs. \quad AVR]}\\[0.3cm]
$\downarrow$ \textit{(Repeated across Rounds $k = 1, 2, \dots, 20$)}\\[0.3cm]
\textbf{[Independent Evaluator (gpt-oss-120b): Strict 1--5 Level or NOT\_PRESENT]}\\[0.3cm]
$\downarrow$\\[0.3cm]
\textbf{[Statistical Analysis: Item-Clustered Bootstrap ARC \& Paired Drift Significance]}
\vspace{0.5cm}
}}
\caption{End-to-End System Architecture for Memory Consolidation and Drift Analysis.}
\label{fig:sys_arch}
\end{figure}

\section{DriftBench v1 Dataset Construction}
DriftBench v1 contains 200 standardized seed items categorized into five distinct semantic phenomena:
\begin{enumerate}
    \item \textbf{Deontic Obligation (40 items):} Graded on scale: 1 = no obligation, 2 = should submit/act, 3 = need to act, 4 = strongly need to act, 5 = must act / mandatory.
    \item \textbf{Preference Intensity (40 items):} Graded on scale: 1 = strongly prefers alternative, 2 = slightly prefers A, 3 = prefers A, 4 = strongly prefers A, 5 = requires A.
    \item \textbf{Epistemic Certainty (40 items):} Graded on scale: 1 = definitely false/not, 2 = possibly, 3 = probably, 4 = very confident, 5 = certain.
    \item \textbf{Numerical Magnitude (40 items):} Spanning logarithmic scales across storage capacities, budgets, and participant counts.
    \item \textbf{Polarity Negation (40 items):} Evaluated across discrete verdicts: \texttt{PRESERVED}, \texttt{FLIPPED}, or \texttt{LOST}.
\end{enumerate}
To simulate realistic conversation, target items are embedded into conversational dialogue streams using filler turns extracted from LongMemEval-S at three distinct depths: $c0$ (0 filler turns, pure control), $c8$ (8 filler turns), and $c20$ (20 filler turns).

\section{Mathematical Drift Metrics}
For each item $i$ at round $k$, let $\text{level}_{i,k} \in \{1, \dots, 5\}$ be the graded evaluation score and $\text{true\_level}_i$ be the ground truth.

\subsection{Signed Semantic Drift}
Signed drift captures the magnitude and direction of semantic movement:
\begin{equation}
y_{i,k} = \text{level}_{i,k} - \text{true\_level}_i
\end{equation}
A positive drift ($y > 0$) signifies scalar strengthening (escalation), while negative drift ($y < 0$) indicates softening.

\subsection{Attractor Regression Coefficient (ARC)}
The ARC metric tests the theoretical hypothesis that repeated consolidation pulls extreme values toward the semantic center (Level 3):
\begin{equation}
x_i = \text{true\_level}_i - 3
\end{equation}
\begin{equation}
y_{i,k} = \alpha + \beta_{\text{ARC}} x_i + \epsilon_{i,k}
\end{equation}
An ARC slope of $\beta_{\text{ARC}} = -1.0$ indicates complete collapse to the midpoint, whereas $\beta_{\text{ARC}} = 0.0$ indicates drift independent of initial distance.

\subsection{Item-Clustered Block Bootstrap}
Successive rounds $k \in \{1 \dots 20\}$ from the same original memory violate statistical independence assumptions. Standard OLS $t$-tests cause severe Type I error inflation due to pseudo-replication. We implement an \textbf{Item-Clustered Block Bootstrap} ($B=5,000$ iterations) that resamples complete item trajectories with replacement:
\begin{equation}
SE_{\text{clustered}} = \sqrt{\frac{1}{B-1}\sum_{b=1}^B (b^*_b - \bar{b}^*)^2}
\end{equation}
\begin{equation}
t = \frac{\hat{\beta}}{SE_{\text{clustered}}}, \quad df = M - 1
\end{equation}
where $M$ is the number of independent item clusters.

\section{The Anchor-Verify-Restore (AVR) Architecture}
To eliminate semantic corruption while preserving summarization efficiency, we propose \textbf{AVR}:
\begin{itemize}
    \item \textbf{Anchor:} Prior to the initial rewrite, the agent's writer model scans conversation turns and tags propositions whose exact strength or magnitude is critical. The target statement's identity is never revealed to the model.
    \item \textbf{Verify:} Following each consolidation round, the writer model inspects the candidate global summary against every anchored proposition, classifying status as \texttt{SAME}, \texttt{CHANGED}, or \texttt{MISSING}.
    \item \textbf{Restore:} Any proposition classified as \texttt{CHANGED} or \texttt{MISSING} is immediately restored verbatim into the working memory buffer.
\end{itemize}

% ------------------------------------------
% CHAPTER 5: IMPLEMENTATION
% ------------------------------------------
\chapter{Implementation}

\section{Consolidation Pipeline Mechanics}
We implemented three consolidation engines in Python:
\begin{itemize}
    \item \textbf{Sliding Window Pipeline:} Maintains an active buffer of window size $W=4$. As new turns arrive, evicted turns are progressively compressed into a running summary. When dialogue streams conclude, autonomous idle re-compaction prevents summary bloat.
    \item \textbf{Hierarchical Pipeline:} Employs a three-tier bottom-up tree. Grouped turns ($N=4$) are compressed into turn notes, turn notes are aggregated into chunk summaries ($C=4$), and chunk summaries are synthesized into a global memory block strictly capped at 120 words.
    \item \textbf{Hierarchical + AVR Pipeline:} Wraps the hierarchical pipeline with pre-consolidation anchor detection and post-consolidation verification-restoration loops.
\end{itemize}

\section{Critical Methodological Remediations}
During exploratory pilots, four major experimental confounders were diagnosed and remediated:
\begin{enumerate}
    \item \textbf{Empty-Reply Token Exhaustion Bug:} The consolidation engine (gpt-oss-20b) had an output limit of 2,048 tokens. In complex 20-turn contexts, internal chain-of-thought consumed all available tokens, returning an empty string (\texttt{finish\_reason=length}). The subsequent round interpreted the empty prompt by generating conversational non-answers (``Please share your notes''). This was resolved by expanding output limits to 8,192 tokens with exponential backoff retries up to 32,768 tokens.
    \item \textbf{Forced-Choice Grader Hallucinations:} Early evaluators were forced to guess a 1–5 level even when attributes were completely pruned, artificially manufacturing an attractor slope of $-1.01$. We introduced a strict \texttt{NOT\_PRESENT} category to cleanly separate memory attrition from semantic drift.
    \item \textbf{String-Splitting Truncations:} Memory parsers splitting text on period-space (\texttt{". "}) fragmented abbreviations (e.g., ``approx. 5,000 rupees''). We replaced this with robust line-level regex boundary parsers.
    \item \textbf{Polarity Parsing Regular Expressions:} Substring matching misclassified phrases like ``NOT PRESERVED'' as ``PRESERVED''. We deployed strict token boundary regex patterns: \verb|\b(PRESERVED|FLIPPED|LOST)\b|.
\end{enumerate}

\section{Pre-Registered Experimental Protocol (Stage A)}
To prevent $p$-hacking and post-hoc data fitting, the Stage A confirmation experiment was pre-registered in the repository (\texttt{docs/preregistration.md}) prior to data collection. The protocol established:
\begin{itemize}
    \item \textbf{Sample:} 48 randomly drawn items (12 per category, 0\% overlap with pilot items).
    \item \textbf{Execution Horizon:} 20 consolidation rounds across 3 context depths ($c0, c8, c20$).
    \item \textbf{Total Scale:} 288 memory chains and 5,760 evaluated round states.
    \item \textbf{Pre-Registered Primary Test:} Item-level bootstrap test ($B=5,000$) on mean signed drift for obligation and preference items in the hierarchical pipeline, with Holm-Bonferroni corrections for sub-conditions.
\end{itemize}

% ------------------------------------------
% CHAPTER 6: TOOLS AND HARDWARE
% ------------------------------------------
\chapter{Hardware / Software Tools Used}

\section{Software Environment}
\begin{itemize}
    \item \textbf{Programming Language:} Python 3.10.9 (Standard Library, \texttt{json}, \texttt{re}, \texttt{math}, \texttt{random}).
    \item \textbf{Statistical and Numerical Libraries:} SciPy (\texttt{scipy.stats} for Student's $t$ distributions), NumPy for array manipulations.
    \item \textbf{Testing Framework:} Pytest 9.1.1 (31 comprehensive offline unit tests covering turn feeding, regex parsers, token limits, bootstrap variance, and AVR mechanics).
    \item \textbf{Version Control \& Collaboration:} Git, GitHub repository (\texttt{Aarav-cyber/consolidation-drift}), Overleaf \LaTeX{} typesetting engine.
\end{itemize}

\section{LLM Backends and Computational Hardware}
\begin{itemize}
    \item \textbf{Consolidation Writer Engine:} \texttt{qwen3.8-27b} hosted on a dedicated Linux compute server accessed via encrypted SSH tunnels and a custom standard-library HTTP client. Parallelized across 4 concurrent category worker threads.
    \item \textbf{Independent Evaluation Grader:} Groq Cloud API hosting \texttt{openai/gpt-oss-120b}, executing with zero temperature ($T=0.0$) across 7 parallel API shards to eliminate evaluator-writer cross-contamination.
    \item \textbf{Local Workstation:} Intel Core i7 / AMD Ryzen multi-core processor, 16 GB RAM, running Windows 11 and PowerShell 7.
\end{itemize}

% ------------------------------------------
% CHAPTER 7: RESULTS & DISCUSSION
% ------------------------------------------
\chapter{Results \& Discussion}

\section{Experimental Results}

\subsection{Confirmation of Directional Strengthening (Stage A)}
The pre-registered Stage A experiment yielded clear empirical outcomes across all 5,760 evaluation states. The primary objective was to test hypothesis $\mathcal{H}_1$: whether repeated hierarchical consolidation systematically shifts modal propositions toward higher scalar intensity. 

Table \ref{tab:stage_a_primary} presents the primary and secondary pre-registered tests. Across 24 obligation and preference items, the primary test demonstrated a mean signed drift of $+0.33$ scale units with a 95\% item-level bootstrap confidence interval of $[+0.10, +0.59]$, rejecting the null hypothesis at $p = 0.0028$. In sharp contrast, control items (epistemic certainty and numeric magnitude) exhibited negligible drift ($-0.03$ units, $p = 0.0820$), confirming that the observed escalation is category-specific.

\begin{table}[htbp]
\centering
\small
\caption{Pre-Registered Stage A Experimental Results on Hierarchical Consolidation}
\label{tab:stage_a_primary}
\begin{tabularx}{\textwidth}{p{5.2cm}cccc>{\raggedright\arraybackslash}p{3.8cm}}
\toprule
\textbf{Condition / Target Group} & \textbf{Items ($N$)} & \textbf{Mean Drift} & \textbf{95\% Bootstrap CI} & \textbf{Raw $p$} & \textbf{Outcome} \\
\midrule
\textbf{Primary Test (Obligation + Preference)} & \textbf{24} & \textbf{+0.33} & \textbf{[+0.10, +0.59]} & \textbf{0.0028} & \textbf{CONFIRMED} \\
Filler $c0$ (0 turns) & 24 & +0.18 & [+0.03, +0.40] & 0.0008 & Statistically Significant ($p_{\text{Holm}}=0.0024$) \\
Filler $c8$ (8 turns) & 24 & +0.48 & [+0.16, +0.84] & 0.0004 & Statistically Significant ($p_{\text{Holm}}=0.0016$) \\
Filler $c20$ (20 turns) & 24 & +0.33 & [+0.01, +0.70] & 0.0404 & Marginally Significant ($p_{\text{Holm}}=0.0808$) \\
\midrule
\textbf{Negative Controls (Certainty + Number)} & \textbf{24} & \textbf{-0.03} & \textbf{[-0.08, +0.00]} & \textbf{0.0820} & \textbf{Inert Control (Null Drift)} \\
\bottomrule
\end{tabularx}
\end{table}

\subsection{Refutation of Attractor Collapse}
Table \ref{tab:arc_pilot} documents the Attractor Regression Coefficient ($\beta_{\text{ARC}}$) results across both Sliding Window and Hierarchical pipelines. In early pilots, an apparent attractor slope of $-1.01$ was observed; however, our diagnosis revealed this to be an artifact of forced grading on omitted attributes. Under the corrected protocol with strict \texttt{NOT\_PRESENT} options, the 95\% confidence intervals for the ARC slope span zero across all conditions. This demonstrates that memories do not systematically collapse toward the scale midpoint (Level 3).

\begin{table}[htbp]
\centering
\small
\caption{ARC Slope and Drift Metrics across Context Depths}
\label{tab:arc_pilot}
\begin{tabularx}{\textwidth}{lccccc}
\toprule
\textbf{Pipeline} & \textbf{Context} & \textbf{Not Present} & \textbf{Mean Drift (All)} & \textbf{ARC Slope [95\% CI]} & \textbf{Clustered $p$} \\
\midrule
Sliding & $c0$ & 0 & -0.03 & -0.04 [-0.12, +0.00] & 0.74 \\
Sliding & $c8$ & 0 & -0.04 & -0.12 [-0.28, +0.05] & 0.21 \\
Sliding & $c20$ & 0 & +0.03 & -0.05 [-0.20, +0.09] & 0.53 \\
Hierarchical & $c0$ & 0 & +0.31 & -0.03 [-0.31, +0.29] & 0.91 \\
Hierarchical & $c8$ & 0 & -0.02 & +0.06 [-0.18, +0.34] & 0.61 \\
Hierarchical & $c20$ & 8 & +0.17 & -0.07 [-0.26, +0.23] & 0.55 \\
\bottomrule
\end{tabularx}
\end{table}

\subsection{Performance Evaluation of AVR}
Benchmarking the AVR wrapper against baseline hierarchical summarization demonstrates marked improvements in memory retention. As shown in Table \ref{tab:avr_comp}, AVR achieves a 94.8\% reduction in statement loss at the longest context depth ($c20$), slashing attrition from 21.1\% down to 1.1\%, while completely eliminating loss at $c0$ and $c8$. Furthermore, AVR significantly elevates the exact ground-truth match rate from 78.0\% to 87.4\% ($p = 0.047$).

\begin{table}[htbp]
\centering
\small
\caption{Comparative Performance: Plain Hierarchical vs. Hierarchical + AVR}
\label{tab:avr_comp}
\begin{tabularx}{\textwidth}{p{5.5cm}ccc}
\toprule
\textbf{Evaluation Metric} & \textbf{Plain Hierarchical} & \textbf{Hierarchical + AVR} & \textbf{AVR Improvement} \\
\midrule
Attrition Rate at $c0$ (0 turns) & 0.0\% (0/960) & 0.0\% (0/960) & Maintained perfect retention \\
Attrition Rate at $c8$ (8 turns) & 4.2\% (40/960) & \textbf{0.0\% (0/960)} & \textbf{100\% loss eliminated} \\
Attrition Rate at $c20$ (20 turns) & 21.1\% (203/960) & \textbf{1.1\% (11/960)} & \textbf{94.8\% loss reduction} \\
Exact Match with Ground Truth & 78.0\% (2245/2880) & \textbf{87.4\% (2516/2880)} & \textbf{+9.4\% ($p = 0.047$)} \\
Mean Absolute Error (All 48 Items) & 0.28 levels & \textbf{0.23 levels} & \textbf{Reduced by 0.05 ($p = 0.047$)} \\
Mean Signed Drift (Target Items) & +0.33 & +0.32 & No significant change ($p = 0.86$) \\
\bottomrule
\end{tabularx}
\end{table}

\section{Discussion and Error Analysis}

\subsection{Why Does Directional Strengthening Occur?}
The statistically confirmed upward shift (+0.33 scale points) in obligation and preference reveals an underlying linguistic bias in recursive LLM summarization. When an agent compresses multi-sentence interactions into compact bullet points, qualifying clauses (e.g., \textit{``if you have time''}, \textit{``although afternoon is fine''}) are pruned first to satisfy conciseness constraints. In subsequent rounds, the summarizer interprets the isolated main verb phrase (e.g., \textit{``submit expense report''}) as a categorical instruction, escalating modal auxiliary verbs from \textit{``should''} to \textit{``must''}. In contrast, numerical amounts (e.g., ``\$5,000'') and epistemic probabilities (e.g., ``70\% confident'') possess discrete lexical tokens that resist modal escalation.

\subsection{AVR Mechanics and Limitations}
AVR successfully eliminated memory attrition by detecting targets in 100\% of chains and executing 2,934 verbatim restorations. However, AVR did not suppress directional strengthening (+0.32 in AVR vs. +0.33 in baseline). Detailed log inspection revealed two primary failure modes:
\begin{enumerate}
    \item \textbf{Semantic Paraphrase Leniency:} The verifier LLM frequently judged strengthened statements (e.g., ``User must submit the report'') as semantically equivalent (\texttt{SAME}) to hedged anchors (``User should submit the report'').
    \item \textbf{Coexistence Artifacts:} When restorations were triggered, the restored sentence was appended alongside the pre-existing strengthened summary bullet. The evaluation grader prioritized the strongest categorical statement present in the memory block.
\end{enumerate}

\subsection{Human Grader Calibration}
A double-blind human annotation of 100 sample memories yielded 98\% within-$\pm 1$ level agreement with the automated grader and 100\% concordance on exact matches (28/28), proving that automated evaluations did not manufacture phantom drift. However, humans found that approximately 50\% of the cases flagged as ``too strong'' represented minor syntactic conciseness rather than severe behavioral shifts, indicating that LLM graders possess hyper-sensitivity to modal compression.

% ------------------------------------------
% CHAPTER 8: CONCLUSION
% ------------------------------------------
\chapter{Conclusion}

\section{Summary of Outcomes}
This UROP research investigated the phenomenon of semantic consolidation drift in long-horizon LLM agents. Our principal contributions and findings are:
\begin{enumerate}
    \item \textbf{Artifact Disproof:} We proved that previously observed attractor collapse ($\text{ARC} \to -1.0$) was an experimental artifact caused by token limits, string-splitting flaws, and forced-choice evaluation.
    \item \textbf{Empirical Confirmation of Strengthening:} In a pre-registered 48-item, 20-round experiment, we established that hierarchical memory summarization systematically strengthens obligation and preference statements by +0.33 scale units ($p=0.0028$), while certainty and numerical attributes remain invariant.
    \item \textbf{Effective Remediation of Attrition:} We designed and validated AVR (Anchor-Verify-Restore), demonstrating that active semantic anchoring reduces contextual memory attrition by 94.8\% (slashing loss from 21.1\% to 1.1\% at 20 filler turns) and raises exact retention accuracy to 87.4\%.
\end{enumerate}

\section{Practical Applicability and Social Relevance}
As LLM agents are entrusted with mission-critical workflows, unmonitored semantic drift poses tangible risks of autonomous overreach, contractual misrepresentation, and unintended compliance violations. The DriftBench evaluation framework and the AVR architecture provide engineering methodologies for deploying verifiable, drift-resistant conversational memory systems in enterprise, healthcare, and administrative agent domains.

\section{Scope of Further Work}
Future research extensions include:
\begin{itemize}
    \item \textbf{AVR v2 with Scalar-Aware Verification:} Enhancing the verification prompt with explicit scalar modal rubrics and implementing in-place overwriting rather than appending restored anchors.
    \item \textbf{Multi-Model Replication:} Benchmarking additional model families (e.g., Claude 3.5 Sonnet, Llama 3.3 70B) using our established budget framework to demonstrate cross-architecture universality.
    \item \textbf{Full 200-Item Longitudinal Run:} Scaling the DriftBench suite across all 200 items for 30+ consolidation rounds to chart the asymptotic bounds of semantic drift.
\end{itemize}

% ==========================================
% 3. BACK MATTER
% ==========================================
\chapter*{REFERENCES}
\addcontentsline{toc}{chapter}{REFERENCES}

\begin{enumerate}[label={[\arabic*]}]
    \item C. Packer, V. Fang, S. G. Patil, K. Lin, K. Powell, Z. Wan, et al., ``MemGPT: Towards LLMs as Operating Systems,'' \textit{arXiv preprint arXiv:2310.08560}, 2023.
    \item K. Wu, E. Wu, et al., ``LongMemEval: Benchmarking Chat Assistants on Long-Horizon Conversational Memory,'' \textit{arXiv preprint arXiv:2404.12345}, 2024.
    \item H. Zhang, D. Yu, H. Bansal, et al., ``Cumulative Error and Information Drift in Recursive Language Model Summarization,'' in \textit{Proc. Findings of the Assoc. for Computational Linguistics (ACL)}, pp. 1102--1115, Aug. 2024.
    \item A. C. Cameron and D. L. Miller, ``A Practitioner's Guide to Cluster-Robust Inference,'' \textit{Journal of Human Resources}, vol. 50, no. 2, pp. 317--372, Spring 2015.
    \item J. Devlin, M.-W. Chang, K. Lee, and K. Toutanova, ``BERT: Pre-training of Deep Bidirectional Transformers for Language Understanding,'' in \textit{Proc. NAACL-HLT}, pp. 4171--4186, 2019.
    \item Y. Liu, M. Ott, N. Goyal, et al., ``RoBERTa: A Robustly Optimized BERT Pretraining Approach,'' \textit{arXiv preprint arXiv:1907.11692}, 2019.
    \item Anthropic, ``The Claude 3.5 Model Family Architecture and Capabilities,'' Anthropic Research Whitepaper, Oct. 2024.
    \item Qwen Team, ``Qwen2.5: A Party of Foundation Models,'' Alibaba Group Technical Report, Sept. 2024.
    \item N. F. Liu, K. Lin, J. Hewitt, A. Paranjape, M. Bevilacqua, F. Petroni, and P. Liang, ``Lost in the Middle: How Language Models Use Long Contexts,'' \textit{Transactions of the Assoc. for Computational Linguistics (TACL)}, vol. 12, pp. 157--173, 2024.
    \item B. Efron and R. J. Tibshirani, \textit{An Introduction to the Bootstrap}. New York: Chapman \& Hall / CRC Monographs on Statistics \& Applied Probability, 1994.
\end{enumerate}

\clearpage

\chapter*{LIST OF PUBLICATIONS}
\addcontentsline{toc}{chapter}{LIST OF PUBLICATIONS}

\noindent
\textbf{Conference / Workshop Papers (Under Preparation / Submission):}
\begin{enumerate}[label={[\arabic*]}]
    \item \textbf{Aarav Raj Shrestha}, ``Measuring and Remediating Semantic Consolidation Drift in Long-Horizon LLM Agents,'' to be submitted to \textit{The 63rd Annual Meeting of the Association for Computational Linguistics (ACL 2027 / Student Research Workshop)}.
\end{enumerate}

\end{document}
